% FigTeaser.tex -- Figure 1 (the teaser): the apparatus on top, the ten repositories below.
% Source: fig_teaser_editable_2026-08-20.pdf, cropped to its ink box.
% MUST stay the first float so it numbers as Figure 1.
\begin{figure*}[t]
  \centering
  \includegraphics[width=\linewidth]{fig_teaser}
  \caption{\textbf{\ai{} at a glance.} \emph{Top:} the lifecycle every task runs through, drawn here
  with discrete graph diffusion as the running example. The repository, the model it starts from and
  the cheap proxy metric are frozen; the agent then has four hours on one B300 to read the code,
  change it, and test each idea against that proxy. What it leaves behind is a source-code patch and
  nothing else --- no weights it trained, no cached state. The patch is applied in a fresh container
  and run from initialization for up to twelve hours, and at most the three most recent checkpoints
  are kept. An evaluator fixed before the first run, with no access to the agent's workspace, scores
  them, and the cell takes the best of the three under the task's direction; for the two tasks that
  do not train, the twelve hours are only a ceiling on the verification stage
  (\S\ref{sec:protocol}). \emph{Bottom:} the ten frozen repositories, chosen so that between them
  they cover ten distinct families of training algorithm, each shown with the final metric its own
  evaluator computes and the asset that metric reads (Table~\ref{tab:tasks}).}
  \label{fig:teaser}
\end{figure*}
